
\chapter*{第二版前言}
\addcontentsline{toc}{chapter}{第二版前言}
\markboth{第二版前言}{第二版前言}

这次修订有两个主要目的：一是纠正第一版中出现的种种错误，二是更新我们对本领域近期研究的论述。本书第一版于 2004 年问世以来，辛几何发展迅速。通过 Heegaard Floer 理论 \srcite{318} 以及新近认识到的\BookRecordIndex{T0208}嵌入接触同调与规范理论之间的联系 \srcite{74,219}，辛几何在低维拓扑中找到了新的应用。已有几部重要著作问世，发展出有力的新思想和技巧：Seidel \srcite{371} 论述 \BookRecordIndex{T0021}Fukaya 范畴，Fukaya--Oh--Ohta--Ono \srcite{128} 论述拉格朗日\BookRecordIndex{T0132}弗洛尔同调，Cieliebak 和 Eliashberg \srcite{63} 则论述复流形与\BookRecordIndex{T0358}辛流形之间的关系。另一项令人振奋的发展是，Tamarkin \srcite{388} 和 Guillermou--Kashiwara--Shapira \srcite{166} 引入了\BookRecordIndex{T0073}层论方法，用以证明辛几何中的基础性刚性结果。在若干具体问题上也取得了重大进展；例如，Taubes \srcite{394} 运用 Seiberg--Witten--Floer 理论解决了三维接触流形的 Weinstein 猜想；Hingston \srcite{176} 和 Ginzburg \srcite{143} 凭借哈密顿动力学与弗洛尔理论的新进展解决了 Conley 猜想；而利用 \BookRecordIndex{T0021}Fukaya 范畴，\BookRecordIndex{T0235}邻近拉格朗日猜想已得到部分解决（贡献者包括 Fukaya--Seidel--Smith \srcite{131} 和 Abouzaid \srcite{1}）。Audin--Damian \srcite{24} 的著作如今提供了哈密顿弗洛尔理论的全面阐述，介绍了在 $c_1=0$ 的流形以及单调情形下建立这一理论所需的全部基本分析。最后，Hofer--Wysocki--Zehnder 的长篇系列论文和著作 \srcite{184,185,186,187,188,189} 发展了研究 \BookRecordIndex{T0003}$J$-全纯曲线理论的一种新的泛函分析方法。他们的工作终将为拉格朗日弗洛尔理论及\BookRecordIndex{T0352}辛场论的各种形式奠定坚实的基础。

我们不会详细讨论这些发展的具体内容。不过，我们已在相关之处更新了各章的导言，并扩充了第 9、11 和 12 章中关于 \BookRecordIndex{T0003}$J$-全纯曲线各种应用的讨论；我们的目标是让读者了解主要的新进展及其中的重要新参与者，而不是作全面综述。

许多修订都相当细微。不过，我们重写了关于\BookRecordIndex{T0099}等周不等式的第 4.4 节、定理 7.2.3 的证明、命题 7.4.8 的证明，以及定理 C.4.2 中\BookRecordIndex{T0129}弗雷德霍姆指标求和公式的证明。我们在第 10 章增加了第 10.9 节，给出一种针对与 $z$ 无关的\BookRecordIndex{T0088}殆复结构的\BookRecordIndex{T0256}粘合定理的新几何表述；在附录 C 扩充了第 C.5 节，加入二维\BookRecordIndex{T0088}殆复结构可积性的证明；在附录 D 扩充了第 D.4 节，将原先第 D.4 和 D.5 节的内容纳入其中，并新增关于零亏格\BookRecordIndex{T0339}稳定曲线\BookRecordIndex{T0246}模空间上同调的第 D.6 节。

我们衷心感谢所有指出第一版错误的读者，尤其感谢 Aleksei Zinger：他寄来了一份格外详尽而有用的意见清单。

\begin{flushright}
Dusa McDuff 与 Dietmar Salamon，2012 年 4 月
\end{flushright}
